\section{Background and Motivation}

\subsection{Workloads in CPU clusters}
% \xiangyu{}}

Modern GPU-centric systems, ranging from single-node dense servers to disaggregated clusters, host a heterogeneous mix of workloads with conflicting resource requirements across the memory, compute, and network hierarchy. 
Generative AI inference has evolved into complex pipelines: Prefill-Decode disaggregation separates compute-intensive prompt processing from memory-bandwidth-bound token generation across networked nodes, requiring precise scheduling to minimize tail latency; KV-cache offloading dynamically spills state to host memory or CXL pools during long-context generation, oscillating between PCIe-bound and HBM-bound phases; and Mixture-of-Experts (MoE) models exhibit sparse, data-dependent activation patterns that demand low-latency paging of expert weights. Neural Network training is equally diverse: while dense distributed training relies on predictable, bulk NCCL/RDMA synchronization, Graph Neural Network (GNN) training and sparse embedding lookups generate highly irregular, pointer-chasing memory access patterns that defy standard hardware prefetchers and thrash TLBs. Vector Search engines further vary based on indexing strategies---HNSW requires low-latency random access similar to graph traversal, while IVF relies on high-throughput scanning. 
\xiangyu{It is OK to roughly mention these scenarios (because people might have already known). However, instead of purely saying what they are doing, we should give one central sentence: they all generate heterogeneous workload. }
In multi-tenant clouds, these workloads compete for shared resources; high-priority inference services demand strict temporal isolation and millisecond-level preemption to meet Service Level Agreements (SLAs), while background batch jobs necessitate aggressive resource oversubscription to maximize total cost of ownership (TCO), creating a tension that static hardware partitioning cannot resolve. 
\xiangyu{The following might need a little bit of rewriting. I do not know the relationship with the previous sentence. Does it mean the previous one needs schedulers not to treat per-GPU as black box? or treat different GPUs differently?} 
At the cluster level, a large body of schedulers such as Gandiva, Tiresias, Optimus and Astraea~\cite{gandiva,tiresias,optimus,astraea} aim to place and time-slice deep-learning jobs across GPUs, but these systems necessarily treat per-GPU resource management as a largely fixed black box.

\subsection{GPU Systems Architecture}

\paragraph{\xiangyu{High-level} User-Space Logic}
High-level frameworks (e.g., PyTorch, TensorFlow) and middleware (e.g., CUDA runtime) operate in user space, managing operation graphs, tensor layouts, and command submission. While they possess rich semantic information---such as the structure of a neural network or the latency target of an inference query---they lack the privileges to enforce these requirements at the hardware level. For instance, a framework might attempt to prioritize a critical stream, but the underlying runtime's opaque heuristics for stream-to-queue mapping can blindly serialize it behind a background job. Similarly, user-space allocators can manage virtual memory pools but cannot control physical page placement or handle page fault interrupts, preventing precise overlapping of data transfer and computation.

\paragraph{The Monolithic Kernel Driver}
The OS driver serves as the
central authority for physical resource arbitration, holding ex-
clusive control over the Memory Management Unit (MMU),
hardware interrupts, and global context scheduling. Despite
this privileged position, modern drivers are typically monolithic and ossified, enforcing static, coarse-grained policies.
For memory management, the driver uses generic algorithms
(e.g., LRU eviction) that ignore workload-specific access patterns.
For scheduling, it employs simple time-slicing (e.g., round-robin) to arbitrate between processes. This “blind’’ management disconnects the driver’s powerful mechanisms from the application’s semantic needs, acting as a primary bottleneck for tail latency and efficiency. 
Prior OS-level work such as TimeGraph and Gdev moved
GPU scheduling and memory control into Linux device drivers~\cite{timegraph,gdev}, and more recent systems based on
kernel-space interception and preemptive scheduling (e.g.,
Zhang et al.~\cite{kernelintercept} and GPREEMPT~\cite{gpreempt})
further refactor driver-level control paths. LithOS goes even
further by re-implementing parts of the GPU runtime stack as
an OS component~\cite{lithos}. However, these systems still
embed fixed policies in C code and do not expose a general,
safe, programmable interface to underlying driver and device
mechanisms.


\paragraph{Device-Side Execution}
Actual computation occurs on the GPU device, managed by a combination of user-written kernels (Device Code) and proprietary firmware/hardware schedulers. Device code defines precise thread-level logic but is typically compiled into static binaries that cannot adapt to runtime conditions like cache thrashing. Meanwhile, the hardware scheduler manages the execution of thousands of threads (warps) but offers no programmable interface for prioritization. This creates a "context silo": the device knows exactly which threads are stalling or consuming bandwidth, but this information is trapped on the chip, inaccessible to the host driver's policy engine.


In summary, modern GPU stacks suffer from tight coupling between resource management mechanisms (e.g.,  page migration, kernel execution lifecycle management, performance event collection) and embedded, hard-coded policy logic within kernel modules, closed-source vendor libraries or GPU device code. This tight coupling severely restricts the ability to adapt policies to diverse and evolving workloads, inhibits rapid experimentation, and fundamentally limits the adaptability of GPU infrastructure in production environments.

\subsection{Limitations of Existing Extensibility Mechanisms}

\xiangyu{I think we can simplify the description of various related works by only referencing them but not mentioning their details.
Probably, it is better to highlight limitations.
}

To mitigate this rigidity, developers currently resort to fragmented extensibility approaches:

\paragraph{User-Space Runtimes and Libraries}
Similar to DPDK for networks, to bypass kernel driver limitations and evolve
rapidly, modern systems like vLLM~\cite{kwon2023efficient} or custom runtimes
implement complex management logic entirely in user space. Fine-grained sharing
frameworks such as Salus~\cite{salus} and UVM-focused runtimes like
PILOT~\cite{pilot} likewise centralize time-sharing and memory placement
decisions above the driver. More recent serving systems like Paella, Pie and
KTransformers explicitly expose programmable interfaces: Paella provides
software-defined GPU scheduling for model serving~\cite{paella}, Pie lets
developers write “inferlets’’ that program the LLM generation loop and KV-cache
management~\cite{pie}, and KTransformers offers a flexible CPU–GPU hybrid
inference runtime with pluggable scheduling and expert-placement
strategies~\cite{ktransformers}. XSched, although positioned as an XPU
scheduling framework, also implements flexible, user-defined policies on top of
its XQueue abstraction~\cite{xsched}. In this sense, many user-space systems
are already programmable frameworks.

% However, all of these approaches share 
\xiangyu{two limitations for these approaches:} First, their programmability is \emph{framework-scoped}:
policies are written against a particular serving system or runtime (e.g.,
Pie, Paella, XSched, KTransformers), and cannot be reused across arbitrary
applications or libraries without porting those applications onto the same
stack. Second, they operate strictly above vendor drivers and firmware and thus
lack a trusted, privileged interface to low-level mechanisms such as page
fault handling, MMU updates, interrupt delivery, and device context
management. As a result, user-space runtimes are forced to defensively
“hoard’’ resources; for example, applications statically pre-allocate GPU
memory at startup to manage their KV-cache or model weights, 
preventing dynamic sharing with other processes, and require additional framework to handle that~\cite{lmcache}. 
% Similarly,
% Hardware partitioning techniques like MIG enforce fixed boundaries that cannot
% adapt to bursty traffic. 
Furthermore, user-space schedulers are confined to
\emph{intra-process} or \emph{intra-framework} decisions; they cannot mitigate
interference from ``noisy neighbors’’ (other tenants or frameworks) contending
for global bandwidth. 
Even within a single process, they are unprivileged:
they cannot intervene at the granularity of hardware interrupts or directly
manipulate physical memory mappings, preventing precise compute-transfer
overlap and true preemption, and may impose substantial performance overhead.
Finally, deploying new policies often necessitates invasive changes to the
serving framework and service restarts. 
\xiangyu{I do not think we should say the following sentence in the motivation section. Comment it for now.}
% In contrast, \sys aims to provide a
% \emph{framework-agnostic}, verifier-backed policy substrate in the GPU driver
% and device execution layer that such user-space systems can target, rather
% than another standalone serving framework.



\paragraph{Binary Instrumentation and Profiling Tools on device}
Tools like NVBit~\cite{nvbit}, Neutrino~\cite{neutrino} or sanitizer frameworks allow for the injection of logic directly into GPU binaries. NVBit
provides dynamic SASS-level binary instrumentation for NVIDIA
GPUs~\cite{nvbit}, while Neutrino attaches programmable probes
at the assembly layer to profile kernels at instruction granularity
with an eBPF-like user experience~\cite{neutrino}. 
\xiangyu{Three limitations for these tools:} 
First of all, while offering high visibility, these tools are primarily designed for offline
profiling rather than online policy enforcement. Second, they incur non-trivial runtime overheads, lack strong safety guarantees (a
bug crashes or hang the application), and typically operate in a read-only manner, unable to actively regulate system behavior such as
memory placement or thread scheduling. Using assembly directly also limited it's expressiveness. Third, recent user-space approaches like eGPU~\cite{egpu} extend eBPF to GPUs for observability but lack privileged driver integration, limiting them to application-level instrumentation and incurring prohibitive overheads (up to $1000\times$, see \S\ref{sec:eval}).

\paragraph{Host-Based Kernel Programmability (CPU eBPF)}
The systems community has successfully used eBPF to customize OS policies on CPUs (e.g., \texttt{sched\_ext}, \texttt{cache\_ext}). However, directly extending this model to GPUs faces two fundamental barriers. First, host-based eBPF treats the GPU as an opaque peripheral; it can trace driver commands but remains blind to the device's internal execution state (e.g., warp divergence), making it impossible to react to on-chip bottlenecks. Second, unlike the Linux CPU scheduler which was refactored to support \texttt{sched\_ext}, monolithic GPU drivers lack safe, programmable extension points to dynamically change it's runtime behavior. Retrofitting mechanism-policy separation into these drivers is non-trivial: resource management logic is tightly coupled with low-level hardware interactions (e.g., MMU invalidation sequences), and simply "opening" the driver risks resource leaks or hardware hangs. A new interface design is required to safely expose these privileged mechanisms without compromising kernel stability.

In summary, these existing methods suffer from fundamental architectural limitations: host-side policies (e.g., memory management, scheduling) lack visibility into GPU-internal runtime states, while device-side policy (e.g., memory sync, thread block management, trace instrumentation) are opaque, vendor-specific, and largely inaccessible to host-side management logic. This fundamental asymmetry in policy expression and execution severely constrains the programmability, observability, and adaptability of GPU resource management.
